\begin{exercise}\label{exer:vi27-01}Show that every nonempty open subset of an integral scheme is dense.\end{exercise}

\begin{exercise}\label{exer:vi27-02}Prove that every nonempty open subscheme of an integral scheme is integral.\end{exercise}

\begin{exercise}\label{exer:vi27-03}For a domain $A$, identify the generic point of $\Spec A$.\end{exercise}

\begin{exercise}\label{exer:vi27-04}Compute the local ring at the generic point of $\Spec k[x,y]$.\end{exercise}

\begin{exercise}\label{exer:vi27-05}Compute $K(\mathbb A^2_k)$.\end{exercise}

\begin{exercise}\label{exer:vi27-06}Compute the function field of $V(y-x^3)\subset\mathbb A^2_k$.\end{exercise}

\begin{exercise}\label{exer:vi27-07}Show that $\Frac(A_f)=\Frac(A)$ for a domain $A$ and $0\ne f\in A$.\end{exercise}

\begin{exercise}\label{exer:vi27-08}Show that $\Spec k[x,y]/(xy)$ is reduced but not integral.\end{exercise}

\begin{exercise}\label{exer:vi27-09}Show that $\Spec k[\varepsilon]/(\varepsilon^2)$ is irreducible but not integral.\end{exercise}

\begin{exercise}\label{exer:vi27-10}Show that a disjoint union of two integral schemes need not be integral.\end{exercise}

\begin{exercise}\label{exer:vi27-11}Prove that $\OO_X(U)\to K(X)$ is injective for nonempty open $U$ in integral $X$.\end{exercise}

\begin{exercise}\label{exer:vi27-12}If two sections on a nonempty open agree on a nonempty subopen, prove they agree everywhere.\end{exercise}

\begin{exercise}\label{exer:vi27-13}For $r=a/b\in\Frac(A)$, identify an open subset on which $r$ is regular.\end{exercise}

\begin{exercise}\label{exer:vi27-14}Prove that the regular locus of a rational function is open.\end{exercise}

\begin{exercise}\label{exer:vi27-15}Embed $A_{\mathfrak p}$ into $\Frac(A)$ for a domain $A$.\end{exercise}

\begin{exercise}\label{exer:vi27-16}Show that a dominant morphism of integral schemes sends generic point to generic point.\end{exercise}

\begin{exercise}\label{exer:vi27-17}Show that a dominant morphism induces an injection of function fields.\end{exercise}

\begin{exercise}\label{exer:vi27-18}For domains $A\to B$, prove the affine dominance criterion using the kernel.\end{exercise}

\begin{exercise}\label{exer:vi27-19}Compute the field map induced by $u\mapsto t^3$ from $\mathbb A^1$ to itself.\end{exercise}

\begin{exercise}\label{exer:vi27-20}Show that a proper closed immersion into an integral scheme is not dominant.\end{exercise}

\begin{exercise}\label{exer:vi27-21}Show that $k[x,y]/(y-x^2)$ has function field $k(x)$.\end{exercise}

\begin{exercise}\label{exer:vi27-22}Show that $k[t^2,t^3]$ and $k[t]$ have the same fraction field.\end{exercise}

\begin{exercise}\label{exer:vi27-23}For an integral finite-type $k$-scheme, show that $K(X)$ is finitely generated as a field extension of $k$.\end{exercise}

\begin{exercise}\label{exer:vi27-24}Show that transcendence degree of $K(X)$ is unchanged by replacing $X$ by a nonempty open subscheme.\end{exercise}

\begin{exercise}\label{exer:vi27-25}Explain why all nonzero rational functions are invertible in $K(X)$.\end{exercise}

\begin{exercise}\label{exer:vi27-26}Give an integral scheme that is not normal and identify its function field.\end{exercise}
